%%%%%%%%%%%% Attribution %%%%%%%%%%%%
% This template was created by 
% Dr. Mahdi A. Khemakhem at PSAU
% and may be copied and used freely for 
% non-commercial purposes.
% 01/15/2023
% This template is created based on
% exam document class created 
% https://math.mit.edu/~psh/exam/examdoc.pdf  
%%%%%%%%%%%%%%%%%%%%%%%%%%%%%%%%%%%%%

\documentclass[addpoints,12pt]{exam}
%%%%%%% Uncomment next line and uncomment the previous one to diplay the exam solution  %%%%%%%
%\documentclass[12pt,addpoints,answers]{exam} % To print with answers
\linespread{1.2} 
\usepackage{gfsartemisia-euler}

\usepackage{amsmath}
\usepackage{graphicx}
\usepackage{setspace}

\usepackage{geometry}
 \geometry{
 a4paper,
 %total={170mm,257mm},
 left=10mm,
 right=10mm,
 top=10mm,
 bottom=25mm,
 }

%%%%%%% To support Arabic language %%%%%%%
\usepackage{arabtex} 
\usepackage{utf8}
\setcode{utf8}

%%%%%%% To support PSAU green color and table patterns %%%%%%%
\usepackage{multirow}
\usepackage[table]{xcolor}	%To can change colour lines and cells in tables
\definecolor{psaugreen}{RGB}{94,162,154}
\usepackage{array}
\newcolumntype{C}[1]{>{\centering\arraybackslash}p{#1}}
\newcolumntype{L}[1]{>{\raggedright\arraybackslash}p{#1}}
\newcolumntype{R}[1]{>{\raggedleft\arraybackslash}p{#1}}

%%%%%%% Define pages footer %%%%%%%
\usepackage{lastpage}
\lfoot{\textit{Final Exam-CS3701-Operation Systems}}
\cfoot{Page \thepage \hspace{1pt} of \pageref{LastPage}}
\rfoot{\textit{Trimester 2, 2022$\slash$2023}}
\runningfootrule

%%%%%%% Command for exam document class %%%%%%%
\newcommand{\tf}[1][{}]{%
\fillin[#1][0.25in]%
}

\begin{document}
%\pagecolor{yellow!5}
%%%%%%% Print marks in left margin %%%%%%%
\pointsinmargin % To print points in margin
\boxedpoints   % or \bracketedpoints

%%%%%%% Include the exam cover defined by the file eaxm_cover.tex %%%%%%%
\include{exam_cover}

\newpage

%%%%%%% Begin exam questions %%%%%%%
\begin{questions}
%%Question 1
\question \textbf{Answer the following} 
    \begin{parts}
        \item Choose the correct answer of the following
        \begin{parts}
            \part[1] If $x = 2$ and $y = 5 $, then $x+y$ \\
            \begin{oneparchoices}
                \choice 2
                \CorrectChoice 7
                \choice 21
                \choice 1
            \end{oneparchoices}
            
            \part[1] If $x = 12$ and $y = 5 $, then $x+y$ \\
            \begin{oneparchoices}
                \CorrectChoice 17 \\
                \choice $\displaystyle 7$ \\ 
                \choice $\displaystyle 21$ \\
                \choice $\displaystyle 1$ 
            \end{oneparchoices}
            
            \part[2] Who invented the pencil
            \begin{oneparcheckboxes}
                \choice Me
                \choice You
                \choice They
                \choice Him
                \CorrectChoice Socrates
            \end{oneparcheckboxes}
        \end{parts}
    
        \part[2] If $f(x) = \sin(x)$, then $f'(x) = $\fillin.
        \part True or False
            \begin{parts} 
                \part[2] \tf[T] The world is all that is the case.
                \part[1] \tf[F] My favorite color is blue.
            \end{parts}
    \end{parts}
%\clearpage

\singlespacing

%%Question 2
\question \textbf{Let $\displaystyle f(x) = \sin(x) + x^2$}
    \begin{parts}
        \part[4] Compute $\displaystyle \frac{df}{dx} $.\\
            \fillwithdottedlines{1.5in}
              %%%% To add solutions, add the following and uncomment the seconed line
              %%%% in the document \documentclass[12pt,addpoints,answers]{exam} %     To print with answers
            \begin{solution}
                $f'(x) = \cos(x) + 2 x $
            \end{solution}
        \part[5] Compute $\displaystyle\int_0^1 f(x) \, dx$.\\
            \fillwithdottedlines{1.5in}
            \begin{solution}
                Put here the answer...
            \end{solution}
    \end{parts}

%%Question 3
\question[10] \textbf{Describe the effect of error propagation on numerical results.}\\
\fillwithdottedlines{2.8in}


%%Question 4
\question
    \begin{parts}
        \part[3] What do you do with $f(x) = x$?
        \fillwithdottedlines{2.8in}
        
        \part[3] Is your answer different if $f(x) = \tan(x)$?
        \fillwithdottedlines{1.9in}
    \end{parts}

%%Question 4
\question[8]
    In no more than one paragraph, explain why the earth is round.
    \makeemptybox{1in}

%%Question 4
\question Answer the following 
    \begin{parts}
        \part[1] \label{eqn1}
            find $f(2.25)$.
            \answerline
        
        \part[1]\label{eqn2}
            Approximate $f'(0)$.
            \answerline
        
        \part[1]\label{eqn3}
            Approximate $\displaystyle \int_{1}^{5} f(x) \: \text{dx}$.
            \answerline

        \part[5] Redo questions \eqref{eqn1}, \eqref{eqn2} and \eqref{eqn3} with $f(1) = 4$ and $f(2) = 5$.
        \fillwithdottedlines{1.2in}
    \end{parts}

\end{questions}

\begin{center}
\fbox{\fbox{\parbox{5.5in}{\centering
Hope you all the best! \\
Dr. X Y Z, \quad \quad
Dr. M N T}}}
\includegraphics[width=0.25\textwidth]{logo}
\end{center}

\end{document}




